\chapter{综合高频面试题}
\newcommand{\qa}[5]{\subsection{#1}\textbf{考察点：}#2\par
\textbf{标准回答：}#3\par\textbf{追问方向：}#4\par
\textbf{常见错误：}#5\par}

\qa{为什么要引入倒格子？}{周期性、傅里叶空间。}
{周期晶格上的函数可按倒格矢展开；衍射条件是动量转移等于倒格矢，Bloch 准动量也只在模倒格矢意义下不同。倒格子把平移对称性变成代数关系。}
{第一布里渊区；晶面间距 $d=2\pi/G$。}{只说“计算方便”，不说明周期相位。}

\qa{Bloch 定理的物理意义是什么？}{平移对称性与准动量。}
{因 $[H,T_{\vb R}]=0$，本征态可同时取为平移本征态，故 $\psi=\ee^{i\vb k\cdot\vb r}u_{\vb k}$。电子受周期势作用，但保留平面波式平移相位，$\hbar\vb k$ 是晶体准动量。}
{$\vb k$ 与 $\vb k+\vb G$；能带指标。}{把 $u_{\vb k}$ 当常数，误称电子完全自由。}

\qa{能隙为什么在布里渊区边界打开？}{Bragg 反射与简并微扰。}
{区边界自由电子态 $|\vb k\rangle,|\vb k-\vb G\rangle$ 简并；周期势 $V_G$ 耦合它们，二阶矩阵本征值为 $E_0+V_0\pm|V_G|$，开隙 $2|V_G|$。}{驻波密度；高维 Bragg 面。}{只说“周期势造成”而不给简并与耦合。}

\qa{声学支和光学支的区别是什么？}{原胞内相位与 $k=0$。}
{长波声学支中原胞内原子同相运动且 $\omega\to0$；光学支反相运动，$k=0$ 仍有恢复力和有限频率。极性晶体的光学模可产生偶极矩并与红外光耦合。}{LO--TO 劈裂；原胞原子数与分支数。}{声称所有光学支都必与可见光强耦合。}

\qa{为什么低温声子热容是 $T^3$？}{态密度与 Bose 分布。}
{三维线性色散给 $g(\omega)\propto\omega^2$；能量积分换元 $x=\hbar\omega/\kb T$ 后 $U\propto T^4$，故 $C_V\propto T^3$，适用 $T\ll\Theta_D$。}{二维/一维指数；Debye 截止。}{把 $T^3$ 归因于“有三个偏振”。}

\qa{为什么电子热容是 $T$？}{费米面薄壳层。}
{Pauli 原理使深层电子无法改变占据，只有距 $E_F$ 约 $\kb T$ 的态参与；参与数 $\propto T$，每个激发能标度 $\propto T$，能量增量 $\propto T^2$，所以热容线性。}{Sommerfeld 系数；有效质量。}{用经典均分定理。}

\qa{Drude 模型为什么失败？}{经典统计与能带缺失。}
{它能给 $\sigma=ne^2\tau/m$ 与简单 Hall 效应，但采用经典粒子和经验碰撞时间，不能解释电子热容、某些正 Hall 系数及复杂温度依赖；Sommerfeld 统计和能带理论分别修正这些缺陷。}{平均自由程；Matthiessen 规则。}{说它“完全错误”。}

\qa{费米面有什么物理意义？}{低能激发边界。}
{它是零温占据态在 $k$ 空间的边界；低温输运、热容和磁响应主要由其附近态及速度、曲率决定，而不只由总电子数决定。}{Luttinger 定理；费米面测量。}{把费米面当真实空间表面。}

\qa{金属、半导体、绝缘体的能带区别是什么？}{填充与带隙。}
{金属费米能穿过能带或能带重叠；后两者零温价带满、导带空且有带隙。半导体带隙较小，热或掺杂载流子可观；绝缘体通常带隙较大。}{Mott 绝缘体；半金属。}{仅用电阻率大小作绝对定义。}

\qa{什么是有效质量？为什么可为负？}{能带曲率与半经典动力学。}
{$1/m^*=\hbar^{-2}\partial_k^2E$，它把晶格作用封装进加速度响应。带顶曲率为负，电子的 $m^*$ 为负；改用正电空穴可等效描述。}{各向异性质量张量。}{把它当电子静质量真的改变。}

\qa{空穴是什么？}{近满带的缺电子激发。}
{移走近满价带中的一个电子后，其余电子的总电荷、能量和晶体动量变化可等效为一个带正电、通常正有效质量的准粒子。}{Hall 符号；电子--空穴复合。}{称为空的空间位置或质子。}

\qa{什么是交换相互作用？}{全同粒子对称性。}
{它源于库仑作用结合费米波函数反对称性；不同自旋排列具有不同空间重叠和能量，可有效写成 $-J\vb S_i\cdot\vb S_j$，并非经典磁偶极力。}{直接交换、超交换。}{断言 $J>0$ 的符号不先说明 Hamiltonian 约定。}

\qa{Curie--Weiss 定律怎么来？}{分子场反馈。}
{独立局域矩给 $M=(C/T)H_{\rm eff}$；平均场取 $H_{\rm eff}=H+\lambda M$，解得 $\chi=C/(T-\lambda C)=C/(T-T_C)$。}{反铁磁 Weiss 温度；平均场局限。}{把 $T_C$ 当无相互作用参数。}

\qa{超导和完美导体有什么区别？}{Meissner 效应。}
{完美导体的零电阻使原有磁通冻结；超导体除零电阻外，在平衡态主动排磁，磁场仅穿透 $\lambda_L$，因此超导是具有相干序参量的新热力学相。}{I/II 型；磁通量子化。}{只用“电阻都为零”回答。}

\qa{Lagrange 方程和 Newton 方程有什么关系？}{等价性与约束处理。}
{D'Alembert 原理将 Newton 方程投影到允许虚位移，理想约束力不作虚功，得到 Lagrange 方程。二者动力学等价，但后者适合广义坐标并显露对称性。}{非保守广义力。}{声称 Lagrange 更“基本”却不说明条件。}

\qa{Hamilton 形式为什么重要？}{正则结构。}
{它把二阶方程化为相空间的一阶正则方程，清晰呈现 Poisson 括号、正则变换和 Liouville 定理，并直接连接统计力学与量子对易关系。}{Hamilton--Jacobi；辛结构。}{总认定 $H=T+V$。}

\qa{Noether 定理说了什么？}{连续对称性与守恒量。}
{作用量每一种连续对称性对应一个守恒流或守恒量；时间平移、空间平移、旋转分别对应能量、动量、角动量。关键是作用量不变而非某条轨迹外观不变。}{规范对称与电荷守恒。}{把离散对称也机械对应连续守恒流。}

\qa{Maxwell 方程每一条的物理意义是什么？}{源、涡旋与守恒。}
{Gauss 电定律说明电荷为电场源；磁 Gauss 定律说明无磁单极；Faraday 定律说明变磁生涡旋电场；Ampère--Maxwell 定律说明电流和变电场生磁场。}{位移电流如何保证电荷守恒。}{遗漏符号方向与适用介质。}

\qa{为什么电磁波是横波？}{真空散度方程。}
{无源均匀真空中平面波代入 $\nabla\cdot\vb E=\nabla\cdot\vb B=0$ 得
$\vb k\cdot\vb E=\vb k\cdot\vb B=0$，旋度方程进一步给三者右手关系。}{波导、等离子体例外。}{把所有介质中的所有模式都称严格横波。}

\qa{规范变换为什么不改变物理场？}{势的描述冗余。}
{代入 $\vb A'=\vb A+\nabla\chi$，旋度中的梯度项为零；代入
$\phi'=\phi-\partial_t\chi$，电场中的两个混合偏导项抵消。}{Lorenz 规范；量子相位。}{把规范变换说成真实场变换。}

\qa{正则系综怎么推导？}{小系统与大热库。}
{总体微正则等概率；热库态数对系统能量作一阶展开，
$\ln\Omega_R(E_{\rm tot}-E_i)=\text{常数}-\beta E_i$，归一化即
$P_i=e^{-\beta E_i}/Z$。}{能量涨落随系统尺度。}{直接假设 Boltzmann 因子而不说明热库。}

\qa{配分函数有什么物理意义？}{微观态加权生成函数。}
{$Z$ 汇总所有微观态的 Boltzmann 权重并完成归一化；$F=-\kb T\ln Z$，对参数求导得到能量、压强、涨落和响应。}{零点选择；Lee--Yang 零点。}{把 $Z$ 直接称作可观测量。}

\qa{Bose 分布和 Fermi 分布有什么区别？}{交换对称性与占据限制。}
{玻色波函数对称、同态占据不限，$\bar n=(e^{\beta(\epsilon-\mu)}-1)^{-1}$；费米波函数反对称、每态最多一粒子，分母为加一。高温低密度时都趋于经典极限。}{BEC；费米能。}{把“每态”误说成每个能级只能一个费米子，忽略简并态。}

\qa{不确定性原理来自哪里？}{不对易与 Hilbert 空间几何。}
{对中心化态 $(A-\langle A\rangle)|\psi\rangle$ 与
$(B-\langle B\rangle)|\psi\rangle$ 用 Cauchy--Schwarz 不等式，得到
$\Delta A\Delta B\ge|\langle[A,B]\rangle|/2$。}{Robertson--Schrödinger 强化式。}{仅解释成测量仪器扰动。}

\qa{微扰论什么时候失效？}{小参数的相对尺度。}
{当微扰矩阵元与相关能级差同量级、存在未处理近简并、耦合引发非解析行为或高阶项不递减时失效；应改用简并微扰、变分或数值对角化。}{渐近级数。}{只比较 $H'$ 与总能量而忽略能级差。}

\qa{简并微扰为什么先对角化？}{零分母与正确基底。}
{简并子空间内 $H_0$ 不选择基底，微扰的一阶作用是其中的矩阵；对角化后才能找到随微扰连续演化的零阶组合并获得一阶劈裂。}{对称性选基。}{任取一组简并本征态直接算对角元。}

\qa{角动量为什么量子化？}{旋转代数与升降算符。}
{对易代数使 $J^2,J_z$ 可共同对角化，$J_\pm$ 每次改变 $m$ 一单位；态范数非负且梯子必须终止，推出 $j$ 为整数或半整数、$m=-j,\ldots,j$。}{SO(3) 与 SU(2)。}{用经典圆周轨道解释自旋量子化。}

\qa{自旋和轨道角动量有什么区别？}{内禀与空间运动。}
{轨道角动量 $\vb L=\vb r\times\vb p$ 来自波函数空间依赖且 $l$ 为整数；自旋是粒子的内禀自由度，可为半整数。两者服从相同角动量代数并可耦合成总角动量。}{自旋轨道耦合。}{把自旋想成有限尺寸小球自转。}

\qa{量子力学中测量的含义是什么？}{Born 规则与状态更新。}
{测量某可观测量得到其谱中的结果，概率由态在相应投影子空间的范数给出；理想选择性测量后态更新到对应子空间。幺正演化与测量更新是公设中的两类规则。}{POVM、退相干。}{把期望值误当单次测量必得值。}

\qa{如何把统计物理和固体物理中的声子联系起来？}{正规模量子化与 Bose 统计。}
{晶格小振动先分解为独立正规模，量子化后每模成为频率 $\omega_{\vb q s}$ 的玻色振子；用 Bose 分布求占据，再乘态密度积分即可求热容和热输运。}{Umklapp 散射；声子化学势为零。}{把声子数当守恒并引入任意化学势。}

\section{口试通用答题模板}
先用一句话给定义或结论；再给最小数学表达；随后明确适用条件；最后补充物理图像和一个局限。遇到追问时优先说明近似层次，而不是继续堆公式。

